% TC-0012 test-case definition (WSP test-case template).
\documentclass[letterpaper]{article}

\usepackage[letterpaper,margin=1in]{geometry}

% Compile from the repository root so this path resolves to the pinned WSP
% submodule.
\input{wsp/testing/test-case-library.tex}

\begin{document}

\def\TCID{TC-0012}
\def\TCTitle{Corpus 217: Animation defaults, Open errors, Play, Stop, Close and events}
\def\TCPurpose{%
  Verify the Animation control through a form: the defaults, Open of a clip, Play with and without arguments, Stop and Close, the errors 53 for a missing file and 5 for a file that is not an AVI file, AutoPlay, Center, BackStyle, and the events Click and DblClick.}
\def\TCPriority{\TCPriorityHigh}

\def\TCRequirementRef{REQ-0005}
\def\TCDesignRef{Not applicable.}
\def\TCFeatureRef{\texttt{tests/corpus/217-mscomctl2-animation.vbp}, \texttt{.frm}, \texttt{.ui}, \texttt{.debug} and \texttt{.expected}.}
\def\TCTestTechnique{%
  Scenario test: a program that prints what it observes, compared with a stored transcript.}

\def\TCPreconditions{%
  A \texttt{wfc} executable that includes this library (WFC version 0.8.28 or later) and the files of the case in \texttt{tests/corpus}.}
\def\TCEnvironment{%
  Windows x64. The \texttt{wfc} executable of WFC, given to CMake as \texttt{WFC\_EXE}; the program is run by the interpreter (\texttt{--interpret}) with the system common controls. The case is registered only when \texttt{WFC\_EXE} is set; without it CTest reports the case as skipped.}
\def\TCAssumptions{%
  Not applicable.}

\def\TCInputData{%
  The clip \texttt{217-clip.avi}, the text file \texttt{217-notavi.txt}, the name \texttt{missing.avi}; the script lines \texttt{click Anim1}, \texttt{dblclick Anim1} and \texttt{click Anim2}.}
\def\TCInitialState{%
  No output of an earlier run is used. CTest test name: \texttt{TC-0012-corpus-animation}; when another project mounts the library the name is prefixed with \texttt{wfc-mscomctl2-}.}

\def\TCProcedure{%
  \begin{enumerate}
    \item Run \texttt{wfc --interpret tests/corpus/217-mscomctl2-animation.vbp}. Form\_Load prints the defaults of both controls, opens the clip, plays it with several argument lists, stops, closes, plays a closed control, opens a missing file and a text file, opens the clip in the AutoPlay control, and sets Center and BackStyle.
    \item The script clicks Anim1, double-clicks Anim1 and clicks Anim2, then closes the form.
    \item Compare the Immediate-window output with \texttt{217-mscomctl2-animation.debug} and standard output with the empty \texttt{217-mscomctl2-animation.expected}.
  \end{enumerate}}
\def\TCExpectedResult{%
  Open of the clip, Play, Stop and Close give error 0, Play on a closed control gives 0, the missing file gives 53 (File not found), the text file gives 5 (Invalid procedure call or argument); the clicks print \texttt{click1}, \texttt{click1}, \texttt{dblclick1} and \texttt{click2}. Both outputs equal the stored files.}
\def\TCPostConditions{%
  No files are written; the interpreter's windows are closed by the script.}

\TCTemplate

\end{document}
